\chapter{Literature Review}
\label{chap:literature_review}

% Writing objective:
% Position the thesis in the literature of vehicle routing, rolling-horizon planning,
% integrated warehouse-routing, and large-neighborhood heuristics.
% Keep the review tightly tied to the thesis method -- do not turn this into a generic survey.

% Main bibliography source:
% - 22.01paper/bibliography.bib
%
% Core papers already curated there include:
% - solomon1987vrptw
% - shaw1998lns
% - ropke2006alns
% - pisinger2010lns
% - pillac2013dvrp
% - archetti2011flexible
% - francis2008pvrp
% - kovacs2014multiperiod
% - coelho2012consistency
% - groer2009consistentvrp
% - subramanian2013richvrp
% - balster2022routepool
% - vangils2018review
% - vangils2021integratedpickroute
% - wouda2024pyvrp

% ============================================================
\section{Classical vehicle routing foundations}
% Write here:
% - Briefly introduce classical VRP / VRPTW foundations.
% - Use Solomon as the anchor for time-window routing.
% - If needed, mention multi-trip behavior as a practical extension, but stay concise.
%
% Must cite:
% - \cite{solomon1987vrptw}
%
% Goal of this section:
% Establish the baseline routing language that later sections extend.

\section{Dynamic, rolling-horizon, and multi-period routing}
% Write here:
% - Explain how the thesis relates to dynamic VRP and rolling planning.
% - Explain why multi-day service flexibility and carryover naturally connect to multi-period routing.
% - If you discuss consistency/stability, connect it to customer-service regularity and planning continuity.
%
% Good citations:
% - \cite{pillac2013dvrp}
% - \cite{francis2008pvrp}
% - \cite{kovacs2014multiperiod}
% - \cite{groer2009consistentvrp}
% - \cite{coelho2012consistency}
%
% Key writing point:
% Your problem is not just dynamic routing; it is rolling-horizon routing with service flexibility and execution-side coupling.

\section{Integrated warehouse-routing and depot-resource coupling}
% Write here:
% - Review literature that combines order picking and transport planning.
% - Explain why this matters for the thesis: picking/gate/staging loads are not external nuisances;
%   they shape routing feasibility and service quality.
%
% Must cite:
% - \cite{vangils2018review}
% - \cite{vangils2021integratedpickroute}
%
% Key writing point:
% This section justifies why the thesis moves beyond route-only optimization.

\section{Heuristic and rich-VRP solution approaches}
% Write here:
% - Review LNS / ALNS / rich-VRP hybrid methods as the main practical solution tradition.
% - Keep the emphasis on methods that support real constraints and scalable heuristic search.
%
% Good citations:
% - \cite{shaw1998lns}
% - \cite{ropke2006alns}
% - \cite{pisinger2010lns}
% - \cite{subramanian2013richvrp}
% - \cite{balster2022routepool}
% - \cite{vidal2013hgs}
% - \cite{wouda2024pyvrp}
%
% Key writing point:
% The thesis belongs more naturally to this heuristic/OR tradition than to black-box end-to-end learning.

\section{Emerging learning-based approaches and why this thesis stays OR-centered}
% Write here:
% - Keep this short.
% - Acknowledge neural combinatorial optimization as a frontier direction.
% - Then explicitly explain why this thesis remains centered on interpretable, controllable,
%   depot-aware heuristic optimization.
%
% Important note:
% Only keep this section if you can make it serve a real methodological contrast.
% Do not let it dominate the chapter.

\section{Literature gap and thesis positioning}
% Write here:
% End the chapter with a clear gap statement.
% Good final message:
% Existing literature gives strong tools for routing, dynamic planning, and integrated picking-routing,
% but the combination of (i) rolling-horizon order admission, (ii) depot time-bucket resource pressure,
% and (iii) integrated cross-layer smoothing remains insufficiently addressed in a thesis-ready operational framework.
%
% This section should directly set up Chapter 3 and Chapter 4.
